\begin{table}[htbp]
\centering
\small
\setlength{\tabcolsep}{5pt}
\caption{Held-out interventions in Pythia-160M. Random controls use equally
sized sets matched by layer; entries are mean fractions or normalized recovery.}
\label{tab:stage-n-interventions}
\begin{tabularx}{\textwidth}{@{}>{\raggedright\arraybackslash}p{0.30\textwidth}Xrr@{}}
\toprule
Intervention & Measured effect & Targeted & Random \\
\midrule
Remove heads favoring old values & Fraction of errors corrected & 0.38 & 0.09 \\
Remove heads favoring the current value & Fraction switching to an old value & 0.04 & 0.03 \\
\addlinespace
Replace inputs to L8H2's key & QK score-gap recovery & 0.86 & 0.04 \\
                           & Answer-score-gap recovery & 0.75 & 0.06 \\
\bottomrule
\end{tabularx}
\vspace{3pt}\parbox{\linewidth}{\footnotesize
Head removal is tested on 73 old-value errors or 81 correct answers,
respectively; key-input replacement uses 53 held-out pairs. Component sets
are selected on separate discovery examples.
Recovery is normalized by the difference between successful and failed runs:
0 means no recovery and 1 means full recovery. Random entries are means, not
upper bounds.}
\end{table}
